\begin{table*}[!t]
  \centering
  \caption{Two-level evaluation standard for RL- and bandit-based FL defenses.}
  \label{tab:padrao}
  \small
  \begin{tabular}{@{}l >{\raggedright\arraybackslash}p{6.3cm} >{\raggedright\arraybackslash}p{4.0cm} >{\raggedright\arraybackslash}p{4.0cm}@{}}
    \toprule
    & Requirement & What it establishes & Status in the corpus ($n=77$) \\
    \midrule

    \multicolumn{4}{@{}l}{\emph{Level 1. Establish the incremental contribution of learning}} \\

    M1 &
    Compare against non-learned policies using the same signal and action space, including a signal-based heuristic and a validation-selected fixed action; declare the baseline set $\mathcal{B}$ and how each baseline was selected or tuned &
    Estimate $\widehat{\Delta}_{\mathrm{L}}(\mathcal{B})$ against informative non-learned controls &
    Six with a same-signal heuristic; nine with a defensive fixed action; none with both \\

    M2 &
    Report $\widehat{\Delta}_{\mathrm{L}}(\mathcal{B},T)$ at multiple predefined evaluation horizons &
    Whether the incremental contribution of learning changes as additional experience becomes available &
    None \\

    M3 &
    Evaluate the learned--non-learned comparison involving the M1 controls across multiple paired seeds and report comparison-level uncertainty using confidence intervals or statistical tests, with multiplicity correction when applicable &
    Whether the observed difference is stable across runs and supported beyond run-to-run variation &
    Among the 15 comparisons involving either M1 control, one uses at least three seeds and reports comparison-level variability; none applies a statistical test to the comparison \\

    M4 &
    Keep test data and malicious-client labels unavailable to reward computation, policy inputs, model selection, and update verification unless explicitly part of the deployment setting &
    Separation between information used by the defense and information reserved for evaluation &
    Thirty-eight fully comply \\

    M5 &
    Report the effective learning budget in algorithm-appropriate units, including transitions, updates or gradient steps, episodes, and exploration schedules as applicable &
    Whether the learning opportunity can be quantified and reproduced &
    Twenty-two fully calculable; 47 partially calculable \\

    \midrule

    \multicolumn{4}{@{}l}{\emph{Level 2. Evaluate the broader properties of the complete defense}} \\

    M6 &
    Ablate signal, memory, and decision granularity separately &
    Which components account for the reported performance &
    Signal: 11; memory: 0; granularity: 0 \\

    M7 &
    Evaluate adaptive attacks, including attacks targeting the learned policy through its state or reward mechanism &
    Robustness under adversaries that adapt to the defense &
    Five evaluate adaptive attacks; none targets the learned policy \\

    M8 &
    Compare against strong static defenses under comparable information and resource assumptions, including access to server-side data when applicable &
    Performance relative to strong alternatives under more comparable conditions &
    Ten include strong baselines from both families considered \\

    M9 &
    Report signal requirements and isolate the computational and operational cost of the learning component &
    Resource requirements and potential privacy exposure introduced by the defense &
    Numerical computational cost reported in 41; agent-training cost rarely isolated \\

    M10 &
    Release code and the configurations needed to reproduce the experiments &
    Independent reproduction and auditing &
    Nine with directly accessible implementations sufficient for reproduction \\

    \bottomrule
  \end{tabular}
\end{table*}
